% Generated draft: requires mathematical and source review.
\documentclass[11pt]{article}
\usepackage[margin=25mm]{geometry}
\usepackage{fontspec}
\setmainfont{Latin Modern Roman}
\usepackage{amsmath,amssymb,mathtools,unicode-math}
\setmathfont{Latin Modern Math}
\usepackage{enumitem,needspace,xurl,hyperref}
\setlength{\parindent}{0pt}
\setlength{\parskip}{5pt}
\setlength{\emergencystretch}{3em}
\newcommand{\N}{\mathbb{N}}
\newcommand{\Z}{\mathbb{Z}}
\newcommand{\Q}{\mathbb{Q}}
\newcommand{\R}{\mathbb{R}}
\newcommand{\C}{\mathbb{C}}
\newcommand{\F}{\mathbb{F}}
\newcommand{\A}{\mathbb{A}}
\newcommand{\PP}{\mathbb{P}}
\newcommand{\E}{\mathbb{E}}
\newcommand{\eps}{\varepsilon}
\newcommand{\dd}{\,\mathrm{d}}
\begin{document}
\section*{UnsolvedMath ID: 30001033}
\textbf{Source problem code:} catalogue entry
\subsection{Definitions and mathematical statement}
Let \(G\) be the group associated in the report with the finite 14-triangle complex, given by
\[
G=\langle x,y\mid r_1,r_2,r_3\rangle,
\]
where the relators are exactly
\[
r_1=yxyxyxy^{-3}x^{-3},
\qquad r_2=yx^{-1}y^{-1}x^{-3}y^2x^{-1}yxy,
\qquad r_3=y^3x^{-1}yxyx^2y^2xyx.
\]
The original record denotes its lower central series by \(\gamma_0(G)=G,\gamma_1(G),\ldots\), ``as defined in the report.'' For every integer \(i\geq2\), the proposed exact index pattern is
\[
[\gamma_i(G):\gamma_{i+1}(G)]=8
\quad\text{if }i\equiv0\text{ or }1\pmod3,
\]
\[
[\gamma_i(G):\gamma_{i+1}(G)]=4
\quad\text{if }i\equiv2\pmod3.
\]
The source calls this equivalent to \(G\) having finite width \(3\).

The indexing convention is material: for an ordinary lower central series indexed from zero, one would set \(\gamma_{i+1}(G)=[\gamma_i(G),G]\), with commutator subgroup generated by \(a^{-1}b^{-1}ab\) for the indicated elements. The extraction does not reproduce the report's recursion, so this is an explicit possible interpretation, not a silent replacement by the common convention \(\gamma_1(G)=G\). A shift of indices would change the residue classes in the conjecture.

The exact definition of width is also absent. Bounded width alone does not generally force this particular periodic pattern; any group-specific equivalence needs the report's context. If the relevant factors are elementary abelian \(2\)-groups and width is measured by their dimensions over \(\F_2\), orders \(8\) and \(4\) correspond to dimensions \(3\) and \(2\) for the stated indices, but the factor structure and treatment of the initial terms are not supplied. No independent converse is asserted.
\subsection{Short English statement}
For the explicitly presented 14-triangle group, do successive series quotients have orders \(8,8,4\) according to the stated index residues for every \(i\geq2\)? Preserve the report's zero-based series convention and verify its finite-width reformulation.
\subsection{Sources}
\begin{itemize}
\item Local UnsolvedMath registry, permanent ID \(30001033\), source code catalogue entry, ``Finite Width of the 14-Triangle Complex Group.'' No individual proposer is supplied.
\item \emph{Geometric Group Theory, Hyperbolic Dynamics and Symplectic Geometry} (2009), Oberwolfach Reports, DOI 10.4171/owr/2008/42, pp. 2337--2406: \url{https://doi.org/10.4171/owr/2008/42}.
\item Supplied reference, ``Cayley graph expanders and groups of finite width'': \url{https://arxiv.org/abs/0809.1560}.
\item Supplied reference, ``Finite-width computations'': \url{https://arxiv.org/abs/1304.4480}.
\item Supplied literature triage, checked 2026-08-21: status ``open.'' It reports that the primary paper still calls finite width three a conjecture and that later computations do not prove the general assertion. This is supplied dated curation, not independently verified; time-sensitive status and cited claims require independent verification.
\end{itemize}
\end{document}
